\documentclass[12pt]{article}
\usepackage{amsmath, amssymb}

\begin{document}

\begin{center}
\textbf{\Large CSE 400: Fundamentals of Probability in Computing}
\end{center}

\section*{Lecture 26: Stochastic Processes (SSS, WSS, Ergodic)}

\section*{1. Basics of Stochastic Processes [Revision]}

\subsection*{1.1 Collection of Waveforms}

\subsubsection*{Stochastic Process Definition}

A stochastic process \( X(\zeta, t) \) is an entire collection of waveforms or functions over time dependent on the outcome of a random experiment.

\begin{itemize}
\item The random output heavily depends on the experiment.
\end{itemize}

\subsubsection*{Example}

Tossing a coin gives heads or tails resulting in specific deterministic waveforms.

\subsection*{1.2 Realizations and Random Variables}

\subsubsection*{Fixed in Time Representation}

\textbf{Step 1:}

A specific measured waveform from the experiment is a \textbf{Realization} of the process.

\textbf{Step 2:}

If we freeze the process at a specific time \( t = t_1 \), the process output evaluates to a single Random Variable:
\[
X(\zeta, t_1) = X(t_1)
\]

\textbf{Step 3:}

By evaluating at multiple time instances \( (t_1, t_2) \), we acquire multiple dependent random variables:
\[
X(t_1), X(t_2)
\]

\textbf{Step 4:}

The dependency between these generated random variables dictates the fundamental nature of the stochastic process.

\subsection*{1.3 Characterization of Stochastic Process}

\subsubsection*{Probability Density Functions (PDF)}

\textbf{Step 1: 1st Order Characterization}
\[
t_1 \rightarrow f_{X_1}(x, t_1) \rightarrow F_{X_1}(x, t_1)
\]

\textbf{Step 2: 2nd Order Characterization}
\[
(t_1, t_2) \rightarrow f_{X_1 X_2}(x_1, x_2, t_1, t_2)
\]

\textbf{Step 3: nth Order Characterization}
\[
(t_1, \dots, t_n) \rightarrow f_{X_1,\dots,X_n}(x_1,\dots,x_n, t_1,\dots,t_n)
\]

\section*{2. Statistical Moments}

\subsection*{2.1 Mean Vector (First Order Moment)}

\textbf{Step 1: Definition}
\[
\mu_X(t) = E[X(t)]
\]

\textbf{Step 2: Computation using PDF}
\[
\mu_X(t) = \int_{-\infty}^{\infty} x \cdot f_X(x,t)\,dx
\]

\textbf{Step 3: Property}
\[
\mu_X(t) \text{ changes as a function of time}
\]

\subsection*{2.2 Autocorrelation Function}

\textbf{Step 1: Definition}
\[
R_{XX}(t_1, t_2) = E[X(t_1) X^*(t_2)]
\]

\textbf{Step 2: Using Joint Density (2nd Order Characterization)}
\[
R_{XX}(t_1,t_2) = \int_{-\infty}^{\infty} \int_{-\infty}^{\infty}
x_1 x_2 \cdot f_{X_1 X_2}(x_1,x_2,t_1,t_2)\,dx_1 dx_2
\]

\subsection*{2.3 Properties of Autocorrelation}

\textbf{Conjugate Symmetry}

\textbf{Step 1: Property}
\[
R_{XX}(t_1,t_2) = R_{XX}^*(t_2,t_1)
\]

\textbf{Step 2: Proof Sketch}
\[
R_{XX}^*(t_2,t_1) = (E[X(t_2)X^*(t_1)])^*
\]
\[
(E[X(t_2)X^*(t_1)])^* = E[X^*(t_2)X(t_1)]
\]
\[
E[X^*(t_2)X(t_1)] = E[X(t_1)X^*(t_2)] = R_{XX}(t_1,t_2)
\]

\textbf{Positive Definiteness}

Autocorrelation function forms a positive definite matrix.

\subsection*{2.4 Covariance Function}

\textbf{Definition}
\[
C_{XX}(t_1,t_2) = E[(X(t_1) - \mu_X(t_1))(X^*(t_2) - \mu_X^*(t_2))]
\]

\textbf{Normalized Form}
\[
\rho_{XX} = \frac{C_{XX}(t_1,t_2)}{\sigma_X \sigma_X}
\]

\subsection*{2.5 Positive Definite Function}

\textbf{Quadratic Form Definition}
\[
\sum_{i=1}^{n} \sum_{j=1}^{n} a_i a_j^* R_{XX}(t_i,t_j) \ge 0
\]

\textbf{Matrix Expansion}
\[
[a_1^* ; a_2^* ; \dots ; a_n^*]
\begin{bmatrix}
R_{11} & R_{12} & \dots \\
R_{21} & R_{22} & \dots \\
\vdots & \vdots & \ddots
\end{bmatrix}
\begin{bmatrix}
a_1 \\ a_2 \\ \vdots
\end{bmatrix}
> 0
\]

\textbf{Equivalent Form}
\[
a^H R a > 0
\]

\subsection*{2.6 Positive Definite Matrices}

\textbf{Conditions}

\begin{enumerate}
\item A symmetric matrix is positive definite if all eigenvalues are positive.
\item A matrix is positive definite if all leading principal minors are positive:
\[
|A_1|>0,\; |A_2|>0,\; \dots,\; |A_n|>0
\]
\end{enumerate}

\textbf{Example (Step-by-Step)}

\[
A =
\begin{bmatrix}
2 & -1 & 0\\
-1 & 2 & -1\\
0 & -1 & 2
\end{bmatrix}
\]

\textbf{Step 1:}
\[
|A_1| = 2 > 0
\]

\textbf{Step 2:}
\[
|A_2| =
\begin{vmatrix}
2 & -1\\
-1 & 2
\end{vmatrix}
= 4 - 1 = 3 > 0
\]

\textbf{Step 3:}
\[
|A_3| = 2(4-1) + 1(-2 - 0)
\]
\[
= 6 - 2 = 4 > 0
\]

Matrix is positive definite.

\section*{3. Stationary Stochastic Process}

\subsection*{3.1 Strict Sense Stationary (SSS)}

\textbf{Step 1: First Order}
\[
f_X(x,t) = f_X(x,t+\tau) = g(x)
\]
\[
\mu_X(t) = \int x g(x)\,dx = \mu_X = \text{constant}
\]

\textbf{Step 2: Second Order}
\[
f_X(x_1,x_2;t_1,t_2) = f_X(x_1,x_2;t_1+\tau,t_2+\tau)
\]

Choose:
\[
\tau = -t_2
\]

\[
f_X(x_1,x_2;t_1,t_2) = f_X(x_1,x_2;t_1 - t_2,0)
\]

\textbf{Step 3: Autocorrelation Derivation}
\[
R_{XX}(t_1,t_2)
= \int\int x_1 x_2 f_X(x_1,x_2;t_1,t_2)\,dx_1dx_2
\]
\[
= \int\int x_1 x_2 f_X(x_1,x_2;t_1-t_2)\,dx_1dx_2
\]
\[
= \phi(t_1 - t_2)
\]

\textbf{Final Result}
\[
R_{XX}(t_1,t_2) = R_{XX}(t_1 - t_2)
\]

\subsection*{3.2 Wide Sense Stationary (WSS)}

SSS requires all orders to be time-invariant.  
WSS is a weaker condition.

\textbf{Condition 1}
\[
E[X(t)] = \mu_X = \text{constant}, \quad \forall t
\]

\textbf{Condition 2}
\[
E[X(t_1)X^*(t_2)] = R_{XX}(\tau), \quad \tau = t_1 - t_2
\]

\subsection*{3.3 Relationship}
\[
\text{SSS} \Rightarrow \text{WSS}, \quad \text{but} \quad \text{WSS} \not\Rightarrow \text{SSS}
\]

\textbf{Gaussian Process Exception}
\[
\text{WSS} \Rightarrow \text{SSS}
\]

\section*{4. Ergodic Process}

\textbf{Time Average of a Random Process}

(Only heading provided in context; no further derivation present in slides.)

\bigskip
\textbf{END OF SCRIBE}

\end{document}